\begin{pastexam}[label={px:gw3-114-ds-02}]{關務三等資訊處理 114 年 第二題}
（申論；25 分）

假設每個運算元都是一位整數，使用堆疊方法，模擬後序式 $25{\times}42{+}{-}6{+}$ 的計算過程。（25 分）
\end{pastexam}

\begin{pastsoln}
\textbf{方法}（資結書第~6.3.2~節）：由左到右讀後序式，
\begin{itemize}[nosep]
\item 讀到運算元：推入堆疊。
\item 讀到二元運算子 $\circ$：彈出兩個數，\textbf{先彈出的是右運算元} $b$、後彈出的是左運算元 $a$，
算 $a \circ b$，把結果推回堆疊。
\item 讀完時堆疊裡應該恰好剩一個數，就是答案。
\end{itemize}
每個運算元都是一位數，所以 $25{\times}42{+}{-}6{+}$ 的符號依序是 $2, 5, \times, 4, 2, +, -, 6, +$。

\textbf{模擬}（堆疊由左到右寫，右邊是頂端）：
\begin{center}
\begin{tabular}{cclll}
\toprule
步 & 讀到 & 動作 & 計算 & 之後的堆疊 \\
\midrule
1 & 2 & 推入 & & 2 \\
2 & 5 & 推入 & & 2, 5 \\
3 & $\times$ & 彈出 $5$、$2$ & $2 \times 5 = 10$ & 10 \\
4 & 4 & 推入 & & 10, 4 \\
5 & 2 & 推入 & & 10, 4, 2 \\
6 & $+$ & 彈出 $2$、$4$ & $4 + 2 = 6$ & 10, 6 \\
7 & $-$ & 彈出 $6$、$10$ & $10 - 6 = 4$ & 4 \\
8 & 6 & 推入 & & 4, 6 \\
9 & $+$ & 彈出 $6$、$4$ & $4 + 6 = 10$ & 10 \\
\bottomrule
\end{tabular}
\end{center}
讀完後堆疊只剩 $10$，\textbf{計算結果是 $10$}。

對照中序式：這個後序式對應 $\bigl((2 \times 5) - (4 + 2)\bigr) + 6 = 10 - 6 + 6 = 10$，相符。 第 7 步要注意
減法的順序：先彈出的 $6$ 是減數，後彈出的 $10$ 是被減數，算成 $6 - 10$ 就錯了。 整個過程每個符號處理一次，
時間和後序式長度成正比，堆疊最多同時放 $3$ 個數（第 5 步）。
\end{pastsoln}
